\chapter{Conclusion and Perspectives}
\label{chap:conclusion}

\section{Research Questions Revisited}
\textbf{RQ1 -- Coverage.} AnsibleGuard represents 38 review criteria: 15 are
fully automated, 18 have partial support, and 5 require manual judgement. Its
76 analysis concepts and 11 profiles support different review decisions while
retaining ansible-lint and GitLeaks as attributed specialist providers.

\textbf{RQ2 -- Reliability.} The v0.0.6 study analysed all 149 frozen
repositories and recorded 21 partial ansible-lint outcomes. In v0.0.7, 206
of 207 frozen repositories yielded analyzable output; 35 provider runs were
partial and one repository timed out. Neither partial evidence nor a timeout
is reported as a clean scan.

\textbf{RQ3 -- Evolution.} Paired runs on fixed populations isolate changes
in detector behaviour. In v0.0.7, required role-argument reasoning suppressed
247 unsupported variable occurrences across the 206 analyzable repositories;
320 additional comparison identities were reclassifications of existing
findings. These changes do not, by themselves, establish precision.

\textbf{RQ4 -- Variable semantics.} Repository-level analysis distinguishes
missing internal definitions, undeclared inputs, and unresolved runtime
context. A single-agent review supported the role-contract correction and
found a source-first false negative. It is diagnostic evidence, not an
independent accuracy benchmark.

\section{Contribution and Next Validation}
The implemented system combines tolerant Ansible-aware parsing, semantic
graphs, original rules, isolated providers, canonical evidence, and versioned
outputs. Its frozen-corpus workflow separates source observations, provider
health, and human labels. The separate 207-repository numerical benchmark
records 2,986,459 observations across 24 metrics and prepares 16,971 review
rows; candidate workload is measured, but threshold quality is not yet known.

Two independent reviewers should label the agreement pilots, resolve
disagreements, and adjudicate the stratified sample before precision or recall
is estimated. Engineering follow-up should investigate the Miarka variable
false negative and the persistent cosmos-ansible timeout. A held-out
repository study and seeded-defect tests can then assess generalisation and
recall beyond the frozen public corpus.
